% Source pins (SHA-256): PifoStatement.lean=fe9475ebfa2c462d27cb0a99781926074e1281396f2c9e8d32d6dfca788d814e; pifo-elegant-v4.md=356ea38146ced437e20200db04e2ce220c73f9ac9fa892b746f783b8a2d49ab8; pifo-nbe.md=5f7efd28922efc5dedad29a375f030ad6e3fe29f96ba94e77ec308f521541f14; pifo-normalization.md=2d287de7924aa30479218da3e0a9a31577ecb97c5a34395c089f7abed98f61e7; pifo-eta-decision.md=50159fe6aaf8273fc1d811f941de4e2246f7ec84e2270fefc1144c0a890e2ec9; pifo-confluence.md=58274d7fa628cbd04ae265a415cd11bce671e950d53a1653c6e23fe5ee1342ec; pifo-noinduction-codex2.md=97eef33723fdd15eff0003d02ce1cd387a1b1dcd86a42f89302421f0fa614933; pifo-noinduction-codex.md=6aea2cf52bbeb98c855054543edc53531d6b458509b27e176f8690bec97e0a25; pifo-noinduction-fable.md=c90151cf07b362c8f49068920f68dba03fb6ca99723354d717e186ecdb068cd2; pifo-verify-all/VERIFICATION.md=c97ed5daa11c97f71a5d0e3867f96ddabe74438f1f598c049050c735a96b1887.
\section{Mechanization and source discipline}
\label{sec:mechanization}

The formal development separates the proposition being decided from every
attempt to prove it.  The frozen file \lean{PifoStatement.lean} contains only
definitions: entries, queues, topology and state, annotated paths, atomic push
and pop, schedulers, runs, the two equivalences, and finally
\lean{statelessFlushImpliesInterleaved : Prop}.  It contains no imports,
theorems, proofs, \lean{sorry}, or proof-specific helper definitions.  Its
SHA-256 digest is
\[
\sourcehash{fe9475ebfa2c462d27cb0a99781926074e1281396f2c9e8d32d6dfca788d814e},
\]
which exactly matches its accompanying \lean{FROZEN.sha256} entry.

This arrangement makes statement drift mechanically visible.  An independent
proof must target the same final proposition over the same definitions; a
change to an operational convention---for example, whether a failed recursive
pop rolls back or whether \(\None\) is observable---changes the frozen digest.
The paper follows that file even where a more abstract presentation might have
hidden these choices.  \autoref{tab:lean-map} records the definition-level
correspondence.

The consolidated clean-slate verification manifest, \lean{VERIFICATION.md},
has SHA-256
\[
\sourcehash{c97ed5daa11c97f71a5d0e3867f96ddabe74438f1f598c049050c735a96b1887}.
\]
Its inventory records entries 1--16, with cleaned and compressed variants
recorded separately as 7b and 8b, and reports 19 gated results.  It records
eleven gated main-theorem artifacts across eight distinct architectures,
credits seven certificate-free main-theorem proofs, and records two executable
deciders.
Development 14, the NbE route, comprises 13,732 lines across 43 files:
41 proof modules under \lean{PifoNbe/} (13,582 lines), the 21-line root
umbrella \lean{PifoNbe.lean}, and the 129-line \lean{PifoNbe/Audit.lean}; its
answer and all 89 audited declarations are standard-gated or axiom-free.

The manifest's \lean{std} gate permits subsets of
\(\{\lean{propext},\lean{Classical.choice},\lean{Quot.sound}\}\); its
\lean{nat} gate additionally permits the documented \lean{native\_decide}
certification mechanism.  Numbered general developments 7, 8, 9, and 12--16
are \lean{std}-gated, while developments 10 and 11 are \lean{nat}-gated
through the documented L3 certificate; variants 7b and 8b are also
\lean{std}-gated.  Independence is tracked per artifact, and development 7 is
explicitly not counted as independent.  No blanket independence or common
axiom-footprint claim is made.

Development~15 adds a 72-module, 25,260-line canonical-form/normalization
route.  All seven gated endpoints report exactly
\(\{\lean{propext},\lean{Classical.choice},\lean{Quot.sound}\}\), including
\lean{interEquiv\_normalForm}, which proves \(S\intereq N(S)\) for valid
\(S\), \lean{structuralDecision\_correct}, \lean{smallFlushTests\_suffice},
and a tenth gated proof of the frozen main theorem.  Here the structural
procedure of \autoref{sec:structural-decision}, \lean{structuralDecision}, is a
plain executable
\lean{Bool} definition; its correctness for valid schedulers is
kernel-checked, and concrete \lean{\#eval} probes return the expected results.

Development~16 adds a 61-module, 16,323-line eta-comparator route.
\lean{etaCompare} is a plain computable definition whose recursive termination
is kernel-checked by Lean; the clean-slate gate checks its soundness,
completeness, equivalence characterization, and the theorem that every witness
word is duplicate-free and has length two or three.  The soundness and
witness-size theorems use only \(\{\lean{propext},\lean{Quot.sound}\}\).
Its trust-level-zero audit passes, and concrete \lean{\#eval} probes execute the
comparator and replay its separating witnesses through the frozen
\lean{run}/\lean{flushOps} semantics for flat and two-child scheduler pairs.
The gated \lean{PifoEta.answer} is the eleventh main-theorem artifact, the eighth
architecture, the seventh certificate-free answer, and the second executable
decider.

The prose proof assembled here was pinned separately.  Its final dual-refereed
source, \lean{pifo-elegant-v4.md}, has SHA-256
\[
\sourcehash{356ea38146ced437e20200db04e2ce220c73f9ac9fa892b746f783b8a2d49ab8}.
\]
The final revision retains the v2 architecture throughout: the argmin proof of
colored congruence, the deterministic diamond, fixed-head chords with counting
linearization, and swap-before-refine.

The mathematical result sources integrated here are pinned independently.  The
canonical-form and decision development, \lean{pifo-normalization.md}, has
SHA-256
\[
\sourcehash{2d287de7924aa30479218da3e0a9a31577ecb97c5a34395c089f7abed98f61e7}
\]
and 669 lines.  The topology-recursive proof,
\lean{pifo-noinduction-codex2.md}, has SHA-256
\[
\sourcehash{97eef33723fdd15eff0003d02ce1cd387a1b1dcd86a42f89302421f0fa614933}
\]
and 583 lines.  The eta-style decision development,
\lean{pifo-eta-decision.md}, has SHA-256
\[
\sourcehash{50159fe6aaf8273fc1d811f941de4e2246f7ec84e2270fefc1144c0a890e2ec9}
\]
and 546 lines.  All three are the supplied dual-refereed SOUND revisions.  The
first scopes the claims in \autoref{sec:normalization-canonical} and
\autoref{sec:deciding-equivalence}; the second scopes those in
\autoref{sec:infinite-alphabets}; and the third scopes the pair-dependent
algorithm, bounds, and certificates in \autoref{sec:eta-decision}.  The eta
source document predates development~16: that development now encodes and
kernel-checks the recursive comparator, including its termination, soundness,
completeness, and witness-size theorem, as described above.

The confluence and bottom-up development, \lean{pifo-confluence.md}, has
SHA-256
\[
\sourcehash{58274d7fa628cbd04ae265a415cd11bce671e950d53a1653c6e23fe5ee1342ec}
\]
and 698 lines.  It is the source for both the nontermination result and the
summary-carrying normalizer in \autoref{sec:confluence-bottom-up}.  The two
additional construction sketches cited in
\remref{rem:alternative-no-induction} are pinned separately.

The first, \lean{pifo-noinduction-codex.md}, has SHA-256
\[
\sourcehash{6aea2cf52bbeb98c855054543edc53531d6b458509b27e176f8690bec97e0a25}.
\]
The second, \lean{pifo-noinduction-fable.md}, has SHA-256
\[
\sourcehash{c90151cf07b362c8f49068920f68dba03fb6ca99723354d717e186ecdb068cd2}.
\]
The direct source for \autoref{sec:nbe}, \lean{pifo-nbe.md}, has SHA-256
\[
\sourcehash{5f7efd28922efc5dedad29a375f030ad6e3fe29f96ba94e77ec308f521541f14}.
\]
The clean-slate manifest's development 14 records the corresponding NbE-route
kernel check and its exact scope.

The correspondence between individual prose arguments and kernel artifacts is
scoped.  For the global extension, \lean{KConstruct.lean}'s
\lean{no\_chain}, \lean{triangle}, \lean{K}, and \lean{K\_iff} check the
fixed-head argument only for flat normal forms: those theorems assume both
topologies are \lean{flatTopo}.  Exact arbitrary-composite coverage is instead
provided by \lean{Ultra2Extend.preorder\_extension} and
\lean{Ultra2Flush.composite\_extension}, via a different walk-normalization
argument.  Likewise, \lean{Ultra2Main.lean}'s \lean{block\_two\_level},
\lean{block\_flush\_eq}, and \lean{reconstruct} formalize the ingredients of
block reconstruction and the analogous side-derived-family route.  The
compiled \lean{Ultra2Main.step} collapses the two side-derived families
separately and then uses cellwise induction; it does not literally formalize
the prose organization around one arbitrarily preselected common family
\((U_D)_{D\in R}\).

\begin{table}[H]
\centering
\small
\begin{tabularx}{\textwidth}{@{}
  >{\raggedright\arraybackslash}p{.34\textwidth}
  >{\raggedright\arraybackslash}p{.31\textwidth}
  >{\raggedright\arraybackslash}X@{}}
\toprule
Paper result & Pinned source & Location in source\\
\midrule
Operational balance, timestamp compression, ghost occurrences, strict keys
  & elegant v4 & \S1\\
Two-type reduction and recovery table
  & elegant v4 & \S1, ``Two-type lemma''\\
Colored-PIFO argmin coupling
  & elegant v4 & \S1, ``Colored-PIFO lemma''\\
Proper top decomposition and flush equations
  & elegant v4 & \S2\\
Cell autonomy and cross-cell comparison recovery
  & elegant v4 & \S3\\
Three-cell diamond
  & elegant v4 & \S3, ``The three-cell diamond''\\
Fixed-head chords and counting linearization
  & elegant v4 & \S3, ``Global extension''\\
Swap, block refinement, and strict-key collapse
  & elegant v4 & \S4\\
Main finite-alphabet theorem
  & frozen statement; elegant v4 & final definition; \S4\\
Canonical tree, flush determination, and complete invariant
  & normalization v3 & \S\S0--1\\
Structural decision, two- and three-push tests, support lower bound, and closing results
  & normalization v3 & \S\S2--3.4\\
Eta comparator, observation theorem, bounds, and certificates
  & eta decision; consolidated verification manifest & \S\S1--7; Development 16\\
Finite meet decomposition, meet promotion, topology recursion, and finite support
  & \lean{pifo-noinduction-codex2.md} & \S\S2--5\\
Rewrite inventory, nontermination, and equational local confluence
  & \lean{pifo-confluence.md} & \S\S1--3\\
Bottom-up summary algebra, smart constructors, reblocking, and costs
  & \lean{pifo-confluence.md} & \S\S4--8\\
Two alternative alphabet-size-induction-free, topology-recursive constructions
  & \lean{pifo-noinduction-codex.md}; \lean{pifo-noinduction-fable.md}
  & construction summaries\\
Normalization by evaluation and common semantic refinement
  & \lean{pifo-nbe.md} (direct); frozen statement; elegant v4; normalization v3
  & \S\S1--10; semantic dependencies\\
Clean-slate builds, theorem identities, independence, line counts, and axiom gates
  & consolidated verification manifest & Developments 1--16\\
Stateful separation
  & consolidated verification manifest
  & Development 2; \lean{Pifo.lean} and \lean{Extended.lean}\\
\bottomrule
\end{tabularx}
\caption{Result-level traceability for the main proof and integrated sources.}
\label{tab:proof-traceability}
\end{table}

\paragraph{Reproducibility boundary.}
Because pdf\TeX{} embeds build metadata, \lean{paper.pdf} is not
byte-reproducible.  Provenance pins therefore attach to \lean{main.tex},
\lean{sections/*.tex}, and the cited sources; the rendered page count is a
build check, not a PDF-digest content pin.

\paragraph{Counterexample audit.}
The stateful artifact uses a small standalone flat-tree model with packet
identities equal to arrival numbers.  In \lean{Pifo.lean}, the theorems
\lean{diverge\_S1} and \lean{diverge\_S2} kernel-check the two outputs of
\autoref{eq:diverging-ops};
\lean{schedulers\_\allowbreak diverge\_\allowbreak interleaved} records their
inequality, all using \lean{decide}.  The theorem
\lean{flush\_agree\_le\_8} checks both drains and the expected pattern through
length eight, also using \lean{decide}; those core checks report only
\lean{propext}.  In the sibling \lean{Extended.lean},
\lean{flush\_agree\_le\_30} is an additional bounded
\lean{native\_decide} check and therefore carries the expected native-reduction
axiom.  The verification manifest explicitly labels that sweep supporting
evidence rather than a proof for all \(n\).
For that reason, we describe the separating trace and bounded checks as
machine-checked, while the universal flush equivalence in
\propref{prop:stateful-separation} rests on the direct calculation in
\autoref{eq:stateful-flush-pattern}.
